\begin{lemma}
\label{lemma-invert-f-sections}
\begin{slogan}
Sections of quasi-coherent sheaves have only meromorphic singularities
at infinity.
\end{slogan}
Let $X$ be a scheme. Let $f \in \Gamma(X, \mathcal{O}_X)$.
Denote $X_f \subset X$ the open where $f$ is invertible, see
Schemes, Lemma \ref{schemes-lemma-f-open}.
If $X$ is quasi-compact and quasi-separated, the canonical map
$$
\Gamma(X, \mathcal{O}_X)_f \longrightarrow \Gamma(X_f, \mathcal{O}_X)
$$
is an isomorphism. Moreover, if $\mathcal{F}$ is a quasi-coherent
sheaf of $\mathcal{O}_X$-modules the map
$$
\Gamma(X, \mathcal{F})_f \longrightarrow \Gamma(X_f, \mathcal{F})
$$
is an isomorphism.
\end{lemma}

\begin{proof}
Write $R = \Gamma(X, \mathcal{O}_X)$. Consider the canonical morphism
$$
\varphi : X \longrightarrow \Spec(R)
$$
of schemes, see
Schemes, Lemma
\ref{schemes-lemma-morphism-into-affine}.
Then the inverse image of the standard open $D(f)$ on the
right hand side is $X_f$ on the left hand side.
Moreover, since $X$ is assumed quasi-compact and quasi-separated
the morphism $\varphi$ is quasi-compact and quasi-separated,
see Schemes, Lemma \ref{schemes-lemma-quasi-compact-affine} and
\ref{schemes-lemma-compose-after-separated}. Hence by
Schemes, Lemma \ref{schemes-lemma-push-forward-quasi-coherent}
we see that $\varphi_*\mathcal{F}$ is quasi-coherent.
Hence we see that $\varphi_*\mathcal{F} = \widetilde M$
with $M = \Gamma(X, \mathcal{F})$ as an $R$-module.
Thus we see that
$$
\Gamma(X_f, \mathcal{F}) =
\Gamma(D(f), \varphi_*\mathcal{F}) =
\Gamma(D(f), \widetilde M) = M_f
$$
which is exactly the content of the lemma. The first displayed isomorphism
of the lemma follows by taking $\mathcal{F} = \mathcal{O}_X$.
\end{proof}
